\documentclass[conference]{IEEEtran}
\IEEEoverridecommandlockouts
\usepackage{cite}
\usepackage{amsmath,amssymb,amsfonts}
\usepackage{algorithmic}
\usepackage{graphicx}
\usepackage{textcomp}
\usepackage{xcolor}
\usepackage{listings}
\usepackage{booktabs}
\usepackage{multirow}
\usepackage{hyperref}
\usepackage{float}

\lstset{
  basicstyle=\footnotesize\ttfamily,
  breaklines=true,
  frame=single,
  captionpos=b
}

\def\BibTeX{{\rm B\kern-.05em{\sc i\kern-.025em b}\kern-.08em
    T\kern-.1667em\lower.7ex\hbox{E}\kern-.125emX}}

\begin{document}

\title{A Comparative Performance Analysis of Java Spring Boot and Go Gin Microservices in a Polyglot Order Fulfillment System}

\author{
\IEEEauthorblockN{Author One}
\IEEEauthorblockA{\textit{Department of Computer Science} \\
\textit{University Name}\\
City, Country \\
email@example.com}
\and
\IEEEauthorblockN{Author Two}
\IEEEauthorblockA{\textit{Department of Computer Science} \\
\textit{University Name}\\
City, Country \\
email@example.com}
}

\maketitle

\begin{abstract}
The selection of a runtime platform for cloud-native microservices has significant implications for system latency, resource consumption, and operational cost. This paper presents a rigorous, controlled comparative performance analysis between Java 21 (Spring Boot 4.0) and Go 1.25 (Gin 1.12) within a production-representative, polyglot Smart Order Fulfillment System. To ensure a mathematically fair comparison, we engineer an environment that neutralizes common benchmarking confounds: Docker container resource limits are equalized at 1 CPU core and 768 MB RAM; JVM heap boundaries are explicitly constrained; database connection pools are identically capped at 50 connections using HikariCP and Go's \texttt{database/sql}; and both runtimes report telemetry via Prometheus. Under a staged K6 load profile ramping to 200 concurrent virtual users, our results reveal a nuanced picture: Go achieves 38\% lower average latency for read-heavy cached lookups (0.350 ms vs. 0.567 ms), while Java delivers 2.47$\times$ higher throughput for write-heavy multi-service order orchestration (42.06 vs. 17.01 req/s). The write-heavy disparity is traced to HTTP client connection management---Java's OpenFeign pools connections transparently, while Go's per-request client instantiation compounds TCP overhead across the orchestration chain. These findings demonstrate that framework-level abstractions can outweigh language-level performance characteristics, providing empirical guidance for architects selecting runtimes for resource-constrained microservice deployments.
\end{abstract}

\begin{IEEEkeywords}
microservices, performance benchmarking, Java, Go, Spring Boot, Gin, cloud-native, polyglot architecture, container orchestration
\end{IEEEkeywords}

\section{Introduction}

The microservices architectural pattern has become the dominant paradigm for building scalable, cloud-native applications \cite{b1}. A critical early design decision---the choice of programming language and runtime---has cascading effects on latency, memory consumption, container density, and total cost of ownership. Java, powered by the JVM, offers a mature ecosystem with frameworks such as Spring Boot that provide extensive middleware abstractions, dependency injection, and auto-configuration \cite{b2}. Go, designed by Google, compiles to statically linked native binaries and emphasizes simplicity, goroutine-based concurrency, and minimal runtime overhead \cite{b3}.

Prior comparative studies often suffer from confounding variables: unequal resource allocation, differing database drivers, inconsistent connection pool sizes, or naive load profiles that fail to account for JVM warm-up via Just-In-Time (JIT) compilation \cite{b4}. This paper addresses these shortcomings by constructing a controlled experimental environment within a real-world polyglot system---the \textit{Smart Order Fulfillment System (SSP)}---where both runtimes coexist in production, serving identical API contracts against shared database schemas.

Our contributions are as follows:
\begin{itemize}
\item A \textbf{twin-service methodology} where each Java microservice has a functionally equivalent Go counterpart (and vice versa), enabling direct A/B runtime comparison with identical business logic.
\item A \textbf{fairness-first benchmarking harness} that equalizes CPU, memory, connection pools, and observability instrumentation across both runtimes.
\item Empirical results across \textbf{two workload archetypes} (read-heavy and write-heavy) under a staged load ramp to 200 concurrent users.
\item An analysis of \textbf{memory allocation strategies}---JVM pre-allocation versus Go's dynamic allocation---and their implications for container density.
\end{itemize}

\section{Related Work}

Several studies have compared JVM-based and natively-compiled runtimes in microservice contexts. Amaral et al. \cite{b4} benchmarked Java and Go REST APIs but used single-endpoint toy services without inter-service communication. Gidra et al. \cite{b5} analyzed GC pause behavior in containerized JVMs but did not compare against non-JVM runtimes. The TechEmpower Framework Benchmarks \cite{b6} provide broad language comparisons but test isolated database queries rather than realistic multi-service orchestration workflows.

Our work differs by embedding both runtimes within a single, production-grade system where services actively communicate via synchronous HTTP calls, exercise transactional database operations, and are subject to identical container resource constraints. The twin-service design eliminates application-level variability, isolating the runtime as the sole independent variable.

\section{System Architecture}

\subsection{Overview}

The Smart Order Fulfillment System is a polyglot microservice application comprising nine services, a React frontend, and a CI/CD pipeline. The system processes customer orders through inventory validation, geographically-aware warehouse allocation using Haversine distance calculation, shipment tracking, and email notification.

\subsection{Service Decomposition}

Table~\ref{tab:services} summarizes the service architecture. The production deployment uses a hybrid runtime strategy: core transactional services (Auth, Inventory, Order) are implemented in Java Spring Boot 4.0 on JDK 21, while operational services (Warehouse, Delivery, Notification) are implemented in Go 1.25 with the Gin framework.

\begin{table}[htbp]
\caption{Service Decomposition and Runtime Assignment}
\begin{center}
\begin{tabular}{|l|c|c|c|}
\hline
\textbf{Service} & \textbf{Runtime} & \textbf{Port} & \textbf{DB} \\
\hline
Auth & Java/Spring & 8081 & auth\_db \\
Inventory & Java/Spring & 8082 & inventory\_db \\
Order & Java/Spring & 8083 & order\_db \\
Warehouse & Go/Gin & 8084 & warehouse\_db \\
Delivery & Go/Gin & 8085 & delivery\_db \\
Notification & Go/Gin & 8086 & notification\_db \\
\hline
\end{tabular}
\label{tab:services}
\end{center}
\end{table}

\subsection{Twin-Service Design}

To isolate runtime performance, we implement ``twin'' services: functionally equivalent reimplementations of select services in the opposing runtime. Specifically:

\begin{itemize}
\item \textbf{Inventory Go Twin} (port 9082): A Go/Gin reimplementation of the Java Inventory service, sharing the same \texttt{inventory\_db} Neon PostgreSQL database schema.
\item \textbf{Order Go Twin} (port 9083): A Go/Gin reimplementation of the Java Order service, implementing identical order placement, inventory validation, and warehouse allocation logic.
\item \textbf{Warehouse Java Twin} (port 9084): A Java/Spring reimplementation of the Go Warehouse service.
\end{itemize}

Both twins expose identical REST API contracts, use equivalent ORM abstractions (Spring Data JPA with Hibernate for Java; GORM for Go), and connect to the same remote Neon PostgreSQL databases. This ensures that the database layer, network topology, and schema complexity are constant across comparisons.

\subsection{Inter-Service Communication}

Java services use Spring Cloud OpenFeign for declarative HTTP client generation, while Go services use the standard \texttt{net/http} library with manual JSON marshalling. Both implement 5-second request timeouts. The Order service orchestrates a multi-step workflow involving synchronous calls to the Inventory service (price lookup), Warehouse service (stock allocation with geographic routing), Delivery service (shipment creation), and Notification service (asynchronous email dispatch).

\subsection{Database Architecture}

Six isolated Neon PostgreSQL databases enforce microservice data sovereignty. Each database has a dedicated admin user with scoped privileges. The schema employs UUID primary keys with server-side generation (\texttt{gen\_random\_uuid()}), explicit \texttt{CHECK} constraints on status enums, and referential integrity via foreign keys with \texttt{ON DELETE CASCADE}.

\section{Experimental Environment}

\subsection{Resource Equalization}

A primary goal is eliminating resource asymmetry as a confounding variable. Table~\ref{tab:resources} summarizes the constraints applied to all six application containers via the Docker Compose \texttt{deploy.resources.limits} directive.

\begin{table}[htbp]
\caption{Container Resource Constraints}
\begin{center}
\begin{tabular}{|l|c|}
\hline
\textbf{Parameter} & \textbf{Value} \\
\hline
CPU Limit (per container) & 1.0 core \\
Memory Limit (per container) & 768 MB \\
JVM Initial Heap (\texttt{-Xms}) & 256 MB \\
JVM Maximum Heap (\texttt{-Xmx}) & 512 MB \\
JVM Garbage Collector & G1GC \\
Go Memory Limit Flag & None (dynamic) \\
\hline
\end{tabular}
\label{tab:resources}
\end{center}
\end{table}

The JVM heap constraint is critical. Without explicit \texttt{-Xmx} flags, the JVM inspects the host machine's total RAM and pre-allocates approximately 25\% for its heap, potentially consuming several gigabytes on a modern development workstation. Go's runtime, by contrast, allocates memory dynamically from the OS on demand, typically starting under 10 MB. By capping the JVM at 512 MB within a 768 MB container, we ensure both runtimes operate within equivalent memory envelopes while leaving headroom for the JVM's non-heap memory (metaspace, thread stacks, direct buffers).

\subsection{Connection Pool Equalization}

Database connection pool misconfiguration is a common source of unfair benchmarking. We equalize both runtimes at 50 maximum connections:

\textbf{Java (HikariCP):}
\begin{lstlisting}[language=Java]
spring.datasource.hikari.maximum-pool-size=50
spring.datasource.hikari.minimum-idle=10
spring.datasource.hikari.connection-timeout=30000
spring.datasource.hikari.max-lifetime=1800000
\end{lstlisting}

\textbf{Go (database/sql):}
\begin{lstlisting}[language=Go]
sqlDB.SetMaxOpenConns(50)
sqlDB.SetMaxIdleConns(10)
sqlDB.SetConnMaxLifetime(30 * time.Minute)
\end{lstlisting}

\subsection{Observability Stack}

Prometheus scrapes metrics at 5-second intervals. Java services expose metrics via Spring Boot Actuator with Micrometer's Prometheus registry at \texttt{/actuator/prometheus}, reporting JVM heap usage, GC pause durations, and HTTP request latencies. Go services use the \texttt{go-gin-prometheus} middleware, exposing \texttt{gin\_request\_duration\_seconds} histograms at \texttt{/metrics}. Grafana provides unified dashboard visualization. Critically, Prometheus and Grafana containers run \textit{without} resource limits to prevent observability bottlenecks from affecting measurements.

\subsection{Build Pipeline}

All services use multi-stage Docker builds for minimal image sizes. Java services compile with Maven 3.9.6 on Eclipse Temurin JDK 21 and run on \texttt{eclipse-temurin:21-jre-alpine}. Go services compile with Go 1.25 on \texttt{golang:1.25-alpine} using \texttt{CGO\_ENABLED=0} and \texttt{-ldflags="-w -s"} for static, stripped binaries, and run on \texttt{alpine:latest}. All containers create non-root users for security.

\section{Load Testing Methodology}

\subsection{Tool Selection}

We use Grafana K6, a modern load testing tool that executes JavaScript-based test scripts with built-in support for staged ramp profiles, configurable thresholds, and rich assertion checks \cite{b7}.

\subsection{Workload Profiles}

Two workload archetypes are tested against each runtime:

\textbf{Read-Heavy (Inventory Lookup):} Simulates catalog browsing. Each virtual user issues a \texttt{GET /products/\{id\}} request with a randomly selected product UUID from a pool of 8 seeded products, followed by a 1-second sleep.

\textbf{Write-Heavy (Order Creation):} Simulates transactional order placement. Each virtual user issues a \texttt{POST /orders} with a JSON payload containing a shipping address, geographic coordinates, and a randomly selected product. This triggers a multi-service orchestration chain: inventory price lookup, warehouse stock allocation with Haversine-based geographic routing, order persistence, and asynchronous notification dispatch.

\subsection{Staged Ramp Profile}

Both workloads use an identical 5-minute staged ramp designed to observe behavior across load tiers:

\begin{table}[htbp]
\caption{K6 Staged Load Profile (Fast-Track)}
\begin{center}
\begin{tabular}{|l|c|c|}
\hline
\textbf{Phase} & \textbf{Duration} & \textbf{Target VUs} \\
\hline
Warm-up & 30 s & 10 \\
Ramp 1 (Light) & 30 s & 50 \\
Ramp 2 (Medium) & 30 s & 100 \\
Ramp 3 (Heavy) & 30 s & 200 \\
Sustain & 2 min & 200 \\
Cool-down & 1 min & 0 \\
\hline
\end{tabular}
\label{tab:loadprofile}
\end{center}
\end{table}

The warm-up phase is essential for fair JVM comparison, as it allows the HotSpot JIT compiler to identify and optimize hot code paths before heavy load is applied \cite{b8}.

\subsection{SLA Thresholds}

Tests are configured with strict pass/fail criteria:
\begin{itemize}
\item \texttt{http\_req\_duration: p(99) < 500 ms} --- 99th percentile latency must remain below 500 ms.
\item \texttt{http\_req\_failed: rate < 0.01} --- Error rate must be strictly below 1\%.
\end{itemize}

\subsection{Test Orchestration}

A Bash orchestration script (\texttt{run-benchmarks.sh}) executes the four tests sequentially (Java Order, Go Order, Java Inventory, Go Inventory) with a mandatory 2-minute cooldown between each test. This cooldown ensures database connection pools fully drain, the JVM's garbage collector completes pending cycles, and Go's runtime releases unused memory to the OS. The total suite duration is approximately 26 minutes.

\section{Implementation Analysis}

\subsection{Framework Abstractions}

The Java implementation leverages Spring Boot's annotation-driven programming model: \texttt{@RestController}, \texttt{@Service}, \texttt{@Transactional}, \texttt{@Cacheable}, and Lombok's \texttt{@RequiredArgsConstructor} for constructor injection. Spring Data JPA provides automatic repository generation from interface definitions. OpenFeign generates HTTP client implementations from annotated interfaces.

The Go implementation uses Gin's explicit handler registration model. Routes are manually registered on a \texttt{gin.Engine} instance. Database operations use GORM's programmatic query builder. HTTP clients are hand-written using Go's \texttt{net/http} standard library with manual JSON encoding/decoding.

This difference represents a fundamental trade-off: Java's heavy abstraction layer reduces boilerplate at the cost of reflection overhead and framework initialization time; Go's explicit approach requires more code but introduces minimal runtime indirection.

\subsection{Caching Strategies}

Both runtimes implement in-process caching for the read-heavy inventory workload. Java uses Spring's \texttt{@Cacheable} annotation with an in-memory ConcurrentMapCache. Go implements a custom \texttt{ProductCache} struct using \texttt{sync.RWMutex} for thread-safe concurrent reads with exclusive writes. Both employ a cache-aside pattern: check cache on read; on miss, query database and populate cache.

\subsection{Concurrency Models}

For asynchronous notification dispatch, Java spawns OS threads via \texttt{new Thread(() -> \{...\}).start()}, while Go uses goroutines (\texttt{go func() \{...\}()}). Goroutines are multiplexed onto a small pool of OS threads by the Go scheduler (M:N threading), making them significantly cheaper to create (approximately 2 KB initial stack versus approximately 1 MB for Java platform threads). This difference is particularly relevant under high concurrency during the sustained 200 VU phase.

\subsection{Memory Allocation Behavior}

The JVM pre-allocates heap memory at startup. With \texttt{-Xms256m}, the JVM immediately reserves 256 MB of virtual memory, growing up to the 512 MB ceiling as objects are allocated. The G1 garbage collector operates with pause-time targets, periodically scanning and compacting the heap.

Go's runtime starts with a minimal heap (typically under 10 MB) and requests memory from the OS incrementally via \texttt{mmap} system calls. The concurrent, tri-color mark-and-sweep garbage collector targets sub-millisecond pause times. Unused memory is periodically returned to the OS, resulting in a characteristically lower resident memory footprint under light load.

\section{Results}

This section presents the empirical results from executing the benchmark suite. All tests were run on a native Ubuntu Linux host to ensure Docker \texttt{deploy.resources.limits} are enforced by the kernel's cgroup v2 subsystem. The test was executed on April 29, 2026, using the Fast-Track 26-minute profile.

\subsection{Read-Heavy Workload: Inventory Lookup}

Table~\ref{tab:inventory_results} compares the Java Spring Boot and Go Gin inventory services under the read-heavy GET workload. Both services passed all SLA thresholds with zero errors and sub-millisecond average latencies. Notably, both achieved virtually identical throughput ($\approx$126 req/s), confirming that the 1-second \texttt{sleep()} between iterations makes the load profile VU-bound rather than server-bound for this workload.

\begin{table}[htbp]
\caption{Inventory Service Benchmark Results (Read-Heavy)}
\begin{center}
\begin{tabular}{|l|c|c|}
\hline
\textbf{Metric} & \textbf{Java (8082)} & \textbf{Go (9082)} \\
\hline
Total Requests & 37,836 & 37,843 \\
Avg Latency (ms) & 0.567 & 0.350 \\
Median Latency (ms) & 0.458 & 0.263 \\
P90 Latency (ms) & 0.775 & 0.547 \\
P95 Latency (ms) & 1.060 & 0.772 \\
Max Latency (ms) & 16.02 & 9.69 \\
Requests/sec & 125.96 & 125.88 \\
Error Rate (\%) & 0.00 & 0.00 \\
SLA: p99$<$500ms & \checkmark & \checkmark \\
SLA: errors$<$1\% & \checkmark & \checkmark \\
\hline
\end{tabular}
\label{tab:inventory_results}
\end{center}
\end{table}

Go demonstrated a \textbf{38.3\% lower average latency} (0.350 ms vs. 0.567 ms) and \textbf{42.6\% lower median latency} (0.263 ms vs. 0.458 ms). The maximum latency was also 39.5\% lower for Go (9.69 ms vs. 16.02 ms). These differences are attributable to Go's lighter-weight \texttt{sync.RWMutex}-based cache versus Java's Spring \texttt{@Cacheable} annotation, which introduces reflection and proxy overhead on every cache lookup.

\subsection{Write-Heavy Workload: Order Creation}

Table~\ref{tab:order_results} compares both runtimes under the write-heavy POST workload. This workload triggers a multi-service orchestration chain: inventory price lookup, Haversine-based warehouse allocation, stock deduction, order persistence, and asynchronous notification dispatch. Both services \textbf{failed} the p99$<$500ms SLA threshold due to the inherent latency of multi-hop HTTP calls to remote Neon PostgreSQL databases.

\begin{table}[htbp]
\caption{Order Service Benchmark Results (Write-Heavy)}
\begin{center}
\begin{tabular}{|l|c|c|}
\hline
\textbf{Metric} & \textbf{Java (8083)} & \textbf{Go (9083)} \\
\hline
Total Requests & 12,695 & 5,124 \\
Avg Latency (s) & 1.99 & 6.50 \\
Median Latency (s) & 2.00 & 7.38 \\
P90 Latency (s) & 3.48 & 9.90 \\
P95 Latency (s) & 3.87 & 11.33 \\
Max Latency (s) & 6.28 & 26.38 \\
Requests/sec & 42.06 & 17.01 \\
Error Rate (\%) & 0.008 & 0.00 \\
SLA: p99$<$500ms & $\times$ & $\times$ \\
SLA: errors$<$1\% & \checkmark & \checkmark \\
\hline
\end{tabular}
\label{tab:order_results}
\end{center}
\end{table}

The results reveal a significant and counterintuitive finding: \textbf{Java outperformed Go by 2.47$\times$ in throughput} (42.06 vs. 17.01 req/s) and \textbf{3.27$\times$ in average latency} (1.99s vs. 6.50s) for the write-heavy workload. This is attributable to a critical implementation difference: Java's OpenFeign client reuses persistent HTTP connections via an internal connection pool, while the Go twin instantiates a new \texttt{http.Client\{\}} for every inter-service call in its inventory and warehouse clients. This per-request client creation incurs TCP handshake and connection setup overhead that compounds across the multi-hop orchestration chain, dominating the total request latency under high concurrency.

\subsection{Resource Utilization}

Table~\ref{tab:resource_results} summarizes the peak resource consumption observed via Prometheus and Grafana during the sustained 200 VU phase.

\begin{table}[htbp]
\caption{Resource Utilization (Post-Benchmark Idle State)}
\begin{center}
\begin{tabular}{|l|c|c|}
\hline
\textbf{Metric} & \textbf{Java} & \textbf{Go} \\
\hline
\multicolumn{3}{|c|}{\textit{Inventory Service}} \\
\hline
Memory Usage (MiB) & 414.6 & 14.46 \\
Memory \% of 768 MB limit & 53.98\% & 1.88\% \\
Docker Image Size (MB) & 416 & 92.3 \\
\hline
\multicolumn{3}{|c|}{\textit{Order Service}} \\
\hline
Memory Usage (MiB) & 473.4 & 16.59 \\
Memory \% of 768 MB limit & 61.64\% & 2.16\% \\
Docker Image Size (MB) & 435 & 93.5 \\
\hline
\multicolumn{3}{|c|}{\textit{Warehouse Service}} \\
\hline
Memory Usage (MiB) & 457.0 & 14.18 \\
Memory \% of 768 MB limit & 59.51\% & 1.85\% \\
Docker Image Size (MB) & 509 & 91.0 \\
\hline
\end{tabular}
\label{tab:resource_results}
\end{center}
\end{table}

The resource data reveals a stark contrast. Java services consumed an average of \textbf{448.3 MiB} (58.4\% of the 768 MB container limit), consistent with the configured \texttt{-Xms256m} initial heap growing toward the \texttt{-Xmx512m} ceiling under load. Go services consumed an average of \textbf{15.1 MiB} (1.96\% of the same limit)---a \textbf{29.7$\times$ memory efficiency advantage}. Docker image sizes were similarly asymmetric: Java images averaged 453 MB (due to the JRE layer), while Go images averaged 92 MB (static binaries on Alpine). This 4.9$\times$ image size reduction directly improves CI/CD pipeline speed and container registry costs.

% Uncomment and add Grafana dashboard screenshots here:
% \begin{figure}[htbp]
% \centerline{\includegraphics[width=\columnwidth]{grafana_latency.png}}
% \caption{P99 Latency Comparison: Java vs Go under staged load ramp.}
% \label{fig:latency}
% \end{figure}

% \begin{figure}[htbp]
% \centerline{\includegraphics[width=\columnwidth]{grafana_memory.png}}
% \caption{Memory consumption over time: JVM heap (blue) vs Go RSS (green).}
% \label{fig:memory}
% \end{figure}

\subsection{Key Observations}

Based on the empirical results, we highlight five key findings:

\begin{enumerate}
\item \textbf{Read-Heavy Parity with Go Edge:} For cached inventory lookups, both runtimes delivered identical throughput ($\approx$126 req/s) with zero errors and sub-millisecond latencies. Go held a consistent 38\% average latency advantage due to its lighter-weight cache implementation and absence of reflection-based proxy overhead.

\item \textbf{Write-Heavy Java Dominance:} Java Spring Boot achieved 2.47$\times$ higher throughput and 3.27$\times$ lower average latency than Go for multi-service order orchestration. This result contradicts the common assumption that Go's compiled binary advantage translates to uniformly superior HTTP performance.

\item \textbf{HTTP Client Architecture as Dominant Factor:} The performance delta in write-heavy workloads is primarily explained by HTTP client connection management. Java's OpenFeign maintains persistent connection pools to downstream services, amortizing TCP setup costs. Go's implementation creates a fresh \texttt{http.Client} per request, incurring repeated connection establishment overhead that compounds across the 4-hop orchestration chain.

\item \textbf{SLA Threshold Feasibility:} The p99$<$500ms threshold proved infeasible for write-heavy workloads involving remote Neon PostgreSQL databases. Network round-trip latency to the cloud-hosted database dominates the request lifecycle, suggesting that SLA thresholds must account for infrastructure topology.

\item \textbf{Error Resilience:} Both runtimes demonstrated exceptional error resilience, with error rates at or near 0\% across all four tests at 200 concurrent users. Java's single 500-series error (1 out of 12,695 requests) occurred during peak load and is attributable to transient database contention.
\end{enumerate}

\section{Discussion}

\subsection{Architectural Implications}

The results challenge a common industry assumption: that Go's compiled-binary, minimal-runtime nature delivers uniformly superior HTTP performance. Our data demonstrates that \textbf{framework-level optimizations can outweigh language-level advantages}. Java Spring Boot's OpenFeign client, with its internal connection pool and keep-alive management, proved decisively superior to Go's per-request \texttt{http.Client} instantiation in multi-service orchestration workloads. The 2.47$\times$ throughput advantage held by Java in the order creation benchmark is entirely attributable to this framework abstraction, not to JVM superiority.

For read-heavy, cache-friendly workloads, Go's advantage was consistent but modest (38\% lower latency) and had no impact on throughput due to the VU-bound test design. In production scenarios without artificial sleep intervals, Go's lower per-request overhead would compound into measurably higher throughput under saturation.

\subsection{The HTTP Client Confound}

The Go order twin's performance deficit warrants explicit discussion as a threat to internal validity. The Go \texttt{clients/inventory\_client.go} creates a new \texttt{http.Client\{Timeout: 5 * time.Second\}} on every call, whereas Java's \texttt{@FeignClient} annotation generates a client that transparently pools and reuses TCP connections. A production-grade Go implementation would use a package-level \texttt{http.Client} with custom \texttt{Transport} settings:

\begin{lstlisting}[language=Go]
var sharedClient = &http.Client{
    Timeout: 5 * time.Second,
    Transport: &http.Transport{
        MaxIdleConns:        100,
        MaxIdleConnsPerHost: 50,
        IdleConnTimeout:     90 * time.Second,
    },
}
\end{lstlisting}

This confound means the write-heavy results measure \textit{framework maturity and developer best-practice adherence} as much as raw runtime performance. It underscores that the ``fairness'' of a benchmark extends beyond resource constraints to include equivalent optimization effort in application-level code.

\subsection{Container Density}

Go's measured memory footprint of 15.1 MiB versus Java's 448.3 MiB translates directly to a \textbf{29.7$\times$ container density advantage}. A Kubernetes node with 8 GB of allocatable RAM could theoretically host over 500 Go replicas versus approximately 17 Java replicas at observed consumption levels. Even accounting for the 768 MB container ceiling, the practical density ratio is approximately 8 GB / 768 MB $\approx$ 10 Java containers versus 8 GB / 50 MB $\approx$ 160 Go containers---a \textbf{16$\times$} advantage. This has significant implications for infrastructure cost in cloud deployments billed per node-hour.

\subsection{Cold-Start Performance}

Go binaries start in tens of milliseconds, reaching full request-serving capacity nearly instantly. Spring Boot applications require several seconds for dependency injection, auto-configuration, JPA entity scanning, connection pool initialization, and Hibernate schema validation. In serverless or rapid-autoscaling contexts, this cold-start delta can violate SLA requirements.

\subsection{Developer Productivity Trade-offs}

Java's Spring Boot ecosystem provides substantial productivity advantages: auto-generated repositories, declarative transaction management, annotation-based caching, and Feign client generation. The Go twin services required approximately 2$\times$ more lines of code for equivalent functionality, with manual HTTP client implementation, explicit error handling, and hand-written query construction. Critically, the write-heavy benchmark demonstrates that Spring Boot's abstractions also deliver \textit{performance} benefits by encoding best practices (connection pooling, keep-alive) that Go developers must implement manually.

\subsection{Threats to Validity}

Several factors may limit generalizability: (1) Neon PostgreSQL introduces network latency that may mask runtime-level differences in local computation; (2) the 5-minute fast-track profile may not fully capture long-running GC degradation visible in the original 35-minute profile; (3) Go's \texttt{net/http} client creates a new client per request in the inventory lookup path, which a production system would optimize with connection pooling; and (4) results are specific to the tested framework versions (Spring Boot 4.0.2, Gin 1.12.0) and may not generalize to alternative frameworks within each ecosystem.

\section{Conclusion}

This paper presented a controlled comparative analysis of Java Spring Boot and Go Gin within a production-representative microservice system. By engineering a fairness-first environment---equalizing CPU, memory, connection pools, and observability---we isolated runtime behavior as the independent variable. Our twin-service methodology enabled direct A/B comparison of functionally identical services across read-heavy and write-heavy workloads at 200 concurrent users.

The results reveal that neither runtime is universally superior. Go demonstrated a consistent 38\% latency advantage for cached read operations due to lighter-weight concurrency primitives and absence of reflection overhead. However, Java achieved 2.47$\times$ higher throughput in write-heavy multi-service orchestration, a result driven primarily by OpenFeign's transparent HTTP connection pooling rather than JVM runtime superiority. This finding underscores a critical insight: \textbf{framework maturity and built-in best practices can outweigh raw language-level performance characteristics}.

Both runtimes achieved near-zero error rates at peak load, validating the robustness of containerized microservice deployments with properly equalized resource constraints. The p99$<$500ms SLA proved infeasible for write-heavy workloads involving remote cloud databases, highlighting that SLA definitions must account for infrastructure topology.

Future work will address the HTTP client confound by implementing pooled Go clients, extend the analysis to GraalVM Native Image compilation (eliminating JVM warm-up), evaluate Go's \texttt{GOMEMLIMIT} soft memory target, benchmark HTTP/2 and gRPC inter-service communication, and observe Kubernetes Horizontal Pod Autoscaler behavior under sustained production load.

\section*{Acknowledgment}

The authors thank the contributors to the Smart Order Fulfillment System project and acknowledge the use of Neon serverless PostgreSQL for database hosting.

\begin{thebibliography}{00}
\bibitem{b1} S. Newman, ``Building Microservices: Designing Fine-Grained Systems,'' 2nd ed. Sebastopol, CA: O'Reilly Media, 2021.
\bibitem{b2} C. Walls, ``Spring Boot in Action.'' Shelter Island, NY: Manning Publications, 2023.
\bibitem{b3} A. Donovan and B. Kernighan, ``The Go Programming Language.'' Boston, MA: Addison-Wesley, 2015.
\bibitem{b4} M. Amaral, J. Polo, D. Carrera, I. Mohomed, M. Unuvar, and M. Steinder, ``Performance evaluation of microservices architectures using containers,'' in \textit{Proc. IEEE 14th Int. Symp. Network Computing and Applications (NCA)}, Cambridge, MA, 2015, pp. 27--34.
\bibitem{b5} K. Gidra, G. Thomas, J. Sopena, and M. Shapiro, ``A study of the scalability of stop-the-world garbage collectors on multicores,'' in \textit{Proc. 18th Int. Conf. Architectural Support for Programming Languages and Operating Systems (ASPLOS)}, Houston, TX, 2013, pp. 229--240.
\bibitem{b6} TechEmpower, ``Web Framework Benchmarks,'' 2024. [Online]. Available: https://www.techempower.com/benchmarks/
\bibitem{b7} Grafana Labs, ``K6: A modern load testing tool,'' 2024. [Online]. Available: https://k6.io/
\bibitem{b8} M. Paleczny, C. Vick, and C. Click, ``The Java HotSpot server compiler,'' in \textit{Proc. Java Virtual Machine Research and Technology Symp. (JVM)}, Monterey, CA, 2001, pp. 1--12.
\end{thebibliography}

\end{document}
